% Generated draft: requires mathematical and source review.
\documentclass[11pt]{article}
\usepackage[margin=25mm]{geometry}
\usepackage{fontspec}
\setmainfont{Latin Modern Roman}
\usepackage{amsmath,amssymb,mathtools,unicode-math}
\setmathfont{Latin Modern Math}
\usepackage{enumitem,needspace,xurl,hyperref}
\setlength{\parindent}{0pt}
\setlength{\parskip}{5pt}
\setlength{\emergencystretch}{3em}
\newcommand{\N}{\mathbb{N}}
\newcommand{\Z}{\mathbb{Z}}
\newcommand{\Q}{\mathbb{Q}}
\newcommand{\R}{\mathbb{R}}
\newcommand{\C}{\mathbb{C}}
\newcommand{\F}{\mathbb{F}}
\newcommand{\A}{\mathbb{A}}
\newcommand{\PP}{\mathbb{P}}
\newcommand{\E}{\mathbb{E}}
\newcommand{\eps}{\varepsilon}
\newcommand{\dd}{\,\mathrm{d}}
\begin{document}
\section*{UnsolvedMath ID: 20001933}
\textbf{Source problem code:} catalogue entry
\subsection{Definitions and mathematical statement}
A unit-weight stationary geodesic net consists of geodesic edges with the balancing condition
\[
 \sum_{e\ni v}u_{e,v}=0
\]
at each vertex \(v\), where \(u_{e,v}\) is the outgoing unit tangent of an incident edge (each incident half-edge is counted). An embedded \(\Theta\)-graph has two distinct vertices joined by three nonconstant geodesic arcs whose interiors are pairwise disjoint. For three unit-weight edges on a surface, balancing means that the three outgoing tangent directions meet at angles \(2\pi/3\) at each vertex.

Question 11.1.1, attributed to [Has-Mor1996], asks: does every smooth Riemannian metric \(g\) on \(S^2\) admit a stationary \(\Theta\)-graph? The source excerpt omits its preceding definition of stationary net, so the embedded unit-weight convention used here follows the supplied aid; any broader immersed or weighted convention requires checking against the original definition.

The original source reports the weaker positive-curvature alternative: every positively curved Riemannian two-sphere has a stationary \(\Theta\)-net, a stationary figure-eight net, or a stationary eyeglasses net. This alternative does not select the \(\Theta\) topology. The appended Section 11.2 calls a stationary net nontrivial when it has a vertex of valency at least three, equivalently in its convention when it is not a disjoint union of simple closed geodesics; this introduces a subsequent topic rather than an added hypothesis of Question 11.1.1.

The supplied special-case claim assumes an effective isometric \(S^1\)-action on \((S^2,g)\), exactly two fixed poles, and free action elsewhere. In meridian coordinates, write
\[
 g=\dd r^2+a(r)^2\dd\theta^2,\qquad 0<r<L,\quad \theta\in\R/(2\pi\Z),
\]
with the circles at \(r=0,L\) collapsed smoothly to the poles. For every phase \(\theta_0\), the three meridians at
\[
 \theta_0,\quad\theta_0+2\pi/3,\quad\theta_0+4\pi/3
\]
are reported to form an embedded unit-weight stationary \(\Theta\)-graph: all three edges have length \(L\), and the outgoing angles at both poles are exactly \(120\) degrees. The distinct graphs in this construction are parametrized by \(S^1/C_3\), where \(C_3\) acts by phase addition of \(2\pi/3\). No curvature-sign assumption is used in this claim.

The aid also supplies the family
\[
 g_\epsilon=\dd r^2+(\sin r+\epsilon\sin^3r)^2\dd\theta^2,
 \qquad 0\le r\le\pi,
\]
with smooth pole collapse. It singles out \(\epsilon=1/4\) as a smooth example with sign-changing Gaussian curvature; every graph above then has total length \(3\pi\). The record does not specify the full parameter range for its family or for curvature sign change, and no such range is asserted here. These symmetry-based examples do not settle the question for arbitrary metrics.
\subsection{Short English statement}
The general question asks for a stationary theta-shaped geodesic network on every Riemannian two-sphere. The supplied aid constructs a continuous family using three equally spaced meridians when the metric has the stated circle symmetry, including a particular metric with curvature of both signs. Its reported open status and special-case proof require independent review.
\subsection{Sources}
\begin{itemize}
\item Local UnsolvedMath registry, permanent ID 20001933, exact source code catalogue entry; canonical \texttt{aim-geometry-notes.json}, \texttt{notes[270]}, tag \texttt{question}; AIM Workshop Problem Lists, Geometry, Geodesics, source item 11.1, Question 11.1.1. Record created and updated 14 August 2026.
\item AIM Geodesics problem list, \url{https://aimath.org/pastworkshops/geodesicsproblems.pdf}. The supplied verification says the official list defines stationary nets before asking this question, but does not reproduce that definition. The PDF was not accessed.
\item Original reference label [Has-Mor1996], identified as the Hass--Morgan question by the supplied aid and cited for the positive-curvature three-topology alternative. No full bibliographic entry is supplied.
\item Supplied research attempt 1: \texttt{partial\_result}, problem status at run \texttt{open}, rigor label \texttt{proved}, novelty confidence \texttt{low}; classification \texttt{PARTIAL-PROGRESS}, registry status \texttt{partially\_solved}. Its candidate is the circle-symmetric special-case theorem. The claim that the general question remains open through August 2026, even under positive curvature, is the aid's dated assessment.
\item Ulam AI / UnsolvedMath Contributors; CC BY 4.0. Machine-generated research aid, not peer reviewed; independently verify before citation.
\end{itemize}
\end{document}
